\section{The proof chain currently completed}
\label{sec:proved-spine}

“Verified” in this section refers to the \AuthoredTargetCount\ authored targets
in the \SnapshotDate\ snapshot. The count indicates project size, not completion
of every parent claim below. Exact premises and EasyCrypt names for each
proposition appear in \cref{sec:source-index}.

\subsection{Three main lines of progress}

The current work is easiest to understand as three connected lines with different
statuses.

\begin{enumerate}
  \item \textbf{Key/transcript line}: from the packed-key prefix down to raw
        \(\mu\) and the challenge suffix, then up to actual API memory/traces;
  \item \textbf{Signature-codec line}: closing canonical signature-prefix and
        HBZ/rANS suffix inverses down to actual pack/unpack procedures;
  \item \textbf{MINCORE KeyGen line}: transporting the algebraic result of the
        actual matrix/finalize snapshot to the paper's \(R_q\)-module product.
\end{enumerate}

\begin{figure}[ht]
\centering
\begin{tikzpicture}[node distance=7mm and 8mm]
  \node[provennode] (kp) {992-byte VK prefix\\of packed SK};
  \node[provennode, right=of kp] (mu) {generated raw\\\(\mu_{64}\mapsto\mu_{32}\)};
  \node[provennode, right=of mu] (chal) {position 64\\challenge suffix};
  \node[provennode, below=of kp] (api) {actual API copy,\\address, frame};
  \node[partialnode, below=of mu] (trace) {stored \(\mu\) in\\completed traces};
  \node[partialnode, below=of chal] (top) {complete challenge\\and API security};
  \draw[flowarrow] (kp) -- (mu);
  \draw[flowarrow] (mu) -- (chal);
  \draw[flowarrow] (kp) -- (api);
  \draw[flowarrow] (api) -- (trace);
  \draw[openarrow] (mu) -- (trace);
  \draw[openarrow] (trace) -- (top);
  \draw[openarrow] (chal) -- (top);
\end{tikzpicture}
\caption{The key/transcript line. Solid arrows denote completed theorems;
dashed arrows denote open lifts.}
\label{fig:key-transcript-spine}
\end{figure}

\begin{figure}[ht]
\centering
\begin{tikzpicture}[node distance=7mm and 7mm]
  \node[provennode] (prefix) {signature prefix\\canonical inverse};
  \node[provennode, right=of prefix] (leaf) {HBZ coefficient\\\(\leftrightarrow\) symbol leaf};
  \node[provennode, right=of leaf] (pure) {table certificate\\and pure rANS trace};
  \node[provennode, below=of pure] (copy) {actual encoder trace\\and suffix copy};
  \node[provennode, left=of copy] (loops) {actual decoder\\exact-trace refinement};
  \node[provennode, left=of loops] (core) {actual harness\\core inverse};
  \node[provennode, below=of core] (hbz) {production full-HBZ\\wrapper inverse};
  \node[provennode, below=of hbz] (witness) {fixed all-zero/all-6\\Hoare success};
  \node[deferrednode, below=of witness] (hcodec) {\(h\) suffix codec\\Week~16 deferred};
  \node[blockednode, below=of hcodec] (full) {full signature codec\\and Sign--Verify correctness};
  \draw[flowarrow] (prefix) -- (leaf);
  \draw[flowarrow] (leaf) -- (pure);
  \draw[flowarrow] (pure) -- (copy);
  \draw[flowarrow] (copy) -- (loops);
  \draw[flowarrow] (loops) -- (core);
  \draw[flowarrow] (core) -- (hbz);
  \draw[flowarrow] (hbz) -- (witness);
  \draw[openarrow] (witness) -- (hcodec);
  \draw[openarrow] (hcodec) -- (full);
\end{tikzpicture}
\caption{The signature-codec line. Week~11 closed the actual encoder's
successful suffix, Week~12 the actual decoder consuming an exact trace,
Week~13 core composition with actual copying in between, Week~14 the production
full-HBZ wrapper boundary, and Week~15 fixed all-6 failure exclusion.
The \(h\) codec was frozen when Week~16 switched to MINCORE.}
\label{fig:codec-spine}
\end{figure}

\subsection{Stage 1: the verification-key quotient of a packed secret key}

At the pure-array level, define
\[
  \operatorname{VKPrefix}_{992}(sk,vk)
  \Longleftrightarrow
  \forall,0\le i<992,\;sk_i=vk_i
\]
\sourcepath{PackedKeyPrefix.ec} proves these basic facts.

\begin{itemize}
  \item The byte-by-byte copy loop invariant is preserved at the next index;
  \item Copying all 992 bytes establishes the prefix relation;
  \item Writes from index 992 onward leave the completed prefix unchanged.
\end{itemize}

These array lemmas lift to the actual extracted procedure.

\begin{verified}[Prefix of the actual packer]
Under the mode~2 constants,
\theoremname{pack_sk_m23_mode2_vk_prefix} shows that when actual generated
\identifier{_pack_sk_m23} returns, the first 992 bytes of its secret-key array
equal the input verification-key array. The subsequent loops writing the
eta-vector and key preserve this prefix.
\end{verified}

\begin{verified}[Internal KeyGen return]
\theoremname{keypair_full_m23_mode2_return_prefix} and
\theoremname{keypair_internal_mode2_return_prefix} lift this relation to actual
generated KeyGen return values.
\[
  \text{KeyGen returns }(vk,sk)
  \Longrightarrow \operatorname{VKPrefix}_{992}(sk,vk).
\]
\end{verified}

Both are \statusproved\ partial-correctness theorems. They do not prove
termination of the KeyGen rejection loop, its complete output distribution,
or equivalence to the paper's KeyGen. Thus \claimid{FC-KG-TOP} remains
\statuspartial.

\subsection{Stage 2: from local arrays to public memory}

Using the array relation returned by KeyGen as actual API input to Sign/Verify
requires four kinds of arrow.

\begin{enumerate}
  \item Exporters copying local KeyGen VK/SK arrays to external addresses;
  \item The importer bringing external SK into Sign's local array;
  \item The importer bringing external VK into Verify's local array;
  \item A canonical bridge showing that integer and \(W64\) addresses
        describe the same regions.
\end{enumerate}

\sourcepath{ApiKeyMemoryBridge.ec} proves mode~2's 992/1408-byte prefixes and
disjoint frames for the actual copy helpers.

\begin{verified}[Sequential exports of actual raw KeyGen]
\theoremname{keypair_raw_api_exports_matching_prefixes} shows that when actual
\identifier{cryptolab_haetae_mode2_keypair_internal} returns, the first 992
bytes of the external SK and VK regions agree under explicit
valid/canonical/disjoint-region contracts.
\end{verified}

This is stronger than merely saying that a copy function copies, but weaker
than complete memory safety. Pointer validity and separation are explicit
premises, and KeyGen termination is still not a conclusion.

\subsection{Stage 3: generated raw \texorpdfstring{\(\mu\)}{mu} for Sign and Verify}

Sign absorbs the local packed-SK prefix and squeezes 64 bytes of \(\mu\).
Verify absorbs the VK region in global memory and squeezes 32 bytes.
The actual generated helper chain performs initialization, key/pre/message
absorption, Keccak permutation, finalization, and squeezing through different
procedure interfaces.

\sourcepath{ExtractedMuHashPrefix.ec} proves leaf relations, and
\sourcepath{ExactMuTopControl.ec} lifts them to the actual top procedures.

\begin{verified}[Generated raw \(\mu\) prefix]
Under explicit mode~2 key-prefix, equal pre/message-region, raw-branch,
and valid/no-wrap premises,
\[
  \mathsf{SignRawMu}\sim\mathsf{VerifyHashMu}
\]
holds with conclusion
\[
  \take_{32}(\mu^{64}_{\mathrm{Sign}})
  =\mu^{32}_{\mathrm{Verify}}
\]
The named theorem is
\theoremname{sign_verify_generated_raw_mu_prefix}.
\end{verified}

This theorem places the two actual generated procedures directly on its
theorem surface. It does not simply identify an abstract wrapper proposition
with actual procedures.
In addition, \theoremname{raw_prelen_shr63_zero} and
\theoremname{raw_prelen_branch_guard} show that the actual bit-level branch
selects the raw path.

\subsection{Stage 4: from whole-memory equality to read-region equality}

Once actual Sign has written a signature, there is no reason for the two
executions' entire global memories to agree. Only the key, pre, and message
regions read by the hash are needed.

\begin{verified}[Region-local generated hash]
\theoremname{sign_verify_generated_raw_mu_prefix_regionwise} allows different
global memories. When the key/pre/message read regions are related pointwise,
actual generated Sign/Verify raw-hash return values share the same 32-byte prefix.
\end{verified}

\begin{verified}[Sign output frame]
\theoremname{sign_raw_api_frames_reused_regions} shows that actual raw Sign
preserves VK/SK/pre/message regions disjoint from its outputs even while writing
1474 signature bytes and 8 siglen bytes.
\end{verified}

Together, these theorems establish the memory-level basis that the input bytes
Verify will read remain after Sign. The final transport directly comparing
stored \(\mu\) in two completed traces is still open, however.
That obligation is separated in \cref{sec:boundaries}.

\subsection{Stage 5: the position~64 challenge suffix}

Suppose Sign and Verify have the same challenge state and position, the
position is 64, and the preceding stage gives equal \(\mu_{32}\).
A theorem shows that running each actual generated absorb helper once leaves
equal states and positions.

\begin{verified}[Local challenge-suffix zero-loss]
\theoremname{generated_raw_mu_to_challenge_suffix_zero_loss} sequentially calls
the actual generated raw \(\mu\) procedure and actual generated 32-byte challenge
absorb helper on each side. Under the stated premises, the final
\((\mathit{state},\mathit{position})\) pairs agree.
\end{verified}

Deterministic loss is therefore zero at the raw/internal generated seam.
This theorem does not include the highbits/LSB prefix, complete challenge,
stored \(\mu\) in an actual accepted API trace, or the context branch.

\subsection{Stage 6: actual caller control and accepted Verify}

Weeks~4--6 placed these leaf results into actual public/raw caller contexts.
The following currently compile separately.

\begin{itemize}
  \item Exact \(\mu\) trace and hash-input binding for actual raw Sign;
  \item An exact trace preserving every early-reject/tail branch of actual raw/cryptolab Verify;
  \item The implication from Verify acceptance to reaching the actual tail and hash call;
  \item Binding showing that accepted Verify used actual VK/pre/message descriptors;
  \item An exact trace sequentially executing actual Sign and then actual Verify;
  \item A frame preserving key/pre/message regions during that sequential execution.
\end{itemize}

\begin{verified}[From acceptance to the tail]
For actual raw/cryptolab Verify,
\[
  \mathit{VerifyResult}=0
  \Longrightarrow \mathit{tail\_reached}
\]
is proved. This is neither the reverse implication nor a statement that an
arbitrary 1474-byte signature reaches the tail.
\end{verified}

\begin{verified}[Actual sequential trace]
\theoremname{raw_sign_then_verify_actual_exact_trace} matches return values and
final memory of actual raw Sign followed by actual raw Verify to sequential
execution of two instrumented traces.
\end{verified}

Nevertheless, \claimid{OBL-DIRECT-OBSERVED-MU} is \statuspartial.
An exact trace exposes observation locations and values, but does not
automatically transport a separate return-value pRHL theorem into equality
of ghost fields in two already completed traces.

\subsection{Stage 7: section--retraction for the actual signature prefix}

\sourcepath{Mode2SignaturePrefixPack.ec} and
\sourcepath{Mode2SignaturePrefixUnpack.ec} expose the exact layout of the
actual generated prefix loops.
\sourcepath{Mode2SignaturePrefixRoundTrip.ec} contains the harness calling
both procedures in their actual order.

\begin{verified}[Mode~2 prefix round trip]
For canonical challenge and canonical signed-low inputs,
\theoremname{pack_unpack_sig_prefix_mode2_roundtrip} establishes
\begin{align*}
  \mathit{decodedChallenge}[0..256)&=\mathit{challenge}[0..256),\\
  \mathit{decodedLow}[0..1024)&=\mathit{low}[0..1024)
\end{align*}
and frames the unused decoded-low tail.
\end{verified}

An all-zero witness for the canonical premises also compiles, so the
proposition is nonvacuous. The result nevertheless covers only the signature's
first 1056 bytes. It excludes the two metadata bytes, 416-byte suffix, padding,
and success of the complete unpacker.

\subsection{Stage 8: the HBZ coefficient--symbol leaf}

HBZ input comprises 1024 coefficients \(h_i\in[-6,7)\).
Prepare computes \(s_i=h_i+6\in\{0,\ldots,12\}\), and apply restores
\(s_i-6\) as a 32-bit word.

\begin{verified}[HBZ leaf inverse]
For a canonical HBZ array,
\theoremname{encode_prepare_decode_apply_mode2_inverse} shows that sequentially
executing actual prepare and actual apply restores the first 1024 coefficients,
preserves the decoder coefficient tail, and makes prepare's \(bad\) output zero.
\end{verified}

This is a leaf-level inverse bypassing the rANS core. It alone does not
establish a round trip for the actual compressed suffix.

\subsection{Stage 9: actual tables and pure rANS mathematics}

Mode~2 HBZ's 13 frequencies sum to 1024, and their cumulative intervals
partition \([0,1024)\).
\sourcepath{Mode2HbzTableCertificate.ec} proves this arithmetic as well as
the meaning of reciprocal/bias values and packed decoder words.
Separate files compile theorems that the actual generated literal arrays
satisfy that certificate.

\begin{verified}[Pure single-step inverse]
For each symbol \(s\) and admissible state \(x\), using \(c_s,f_s\) from the
actual table gives
\[
  D_s(E_s(x))=x
\]
The reciprocal-based fast encoder step also equals mathematical \(E_s\)
in the stated normalized range.
\end{verified}

Week~9 defined pure \(xs/cuts/bytes\) traces including normalization bytes
between individual steps.

\begin{verified}[Pure byte stack]
For canonical symbol lists, the following compile.
\begin{itemize}
  \item Normalization emits only 0, 1, or 2 bytes at each step;
  \item Appending emitted segments in decoder order restores the original state;
  \item The 4-byte little-endian initial state is recovered after parsing;
  \item Every trace state stays in \([2^{23},2^{31})\);
  \item The trace's head symbol, reduced state, and normalization segment
        decode one step at a time.
\end{itemize}
\end{verified}

\begin{verified}[Actual copying and control harness]
\theoremname{copy_encoded_suffix_correct} proves pointwise slice copying and
the tail frame of actual \identifier{__copy_encoded_suffix}.
\theoremname{actual_rans_harness_branches_on_encoder_result} shows that the
actual encoder, copy, and decoder are called in that order; the decoder runs
exactly when the encoder's actual \(bad\) result is zero, receiving
\((count,size,m)=(1024,\mathit{size},13)\).
\end{verified}

The control theorem itself concludes neither decoder success nor symbol
equality, final state, or byte-consumption count. Week~12's separate decoder
theorem proves those semantics under the exact-trace input relation, but at
that stage had not yet composed them into the control harness.
Week~13's \theoremname{actual_rans_encode_copy_decode_inverse} adds those
semantic conclusions to the same harness.

\subsection{Stage 10: local semantics of the actual encoder}

Week~10 incorporated \WeekTenTargetCount\ rANS encoder targets into the authored
target manifest and freshly compiled them. Rather than immediately assert one
large theorem between pure traces and actual generated procedures, these
separate array/list correspondence, word arithmetic, inner normalization,
and final serialization \cite{week10-report,rans-encoder-invariant}.

\begin{verified}[Array, word, and serialization lemmas]
\theoremname{encode_trace_suffix_extension} transports the actual 1024-byte
symbol-array suffix to a pure trace recurrence extending it one symbol at a
time. \theoremname{actual_mode2_encoder_word_step_correct} shows that the actual
word expression using literal-table reciprocals, shifts, biases, and complements
equals the pure fast step.
\theoremname{actual_encoder_final_state_serialization} shows that four
\identifier{set8} operations equal little-endian serialization of the 32-bit state.
\end{verified}

\begin{verified}[Actual inner loop and successful-output size]
\theoremname{actual_rans_encode_inner_no_underflow} installs byte-segment and
prefix-frame invariants in the nested normalization loop of actual generated
\identifier{_rans_encode}, establishing no underflow of the \(W64\) cursor.
\theoremname{actual_rans_encode_success_size_bound} is an actual
partial-correctness theorem: if that procedure returns, then
\[
  bad\ne0\quad\lor\quad
  \bigl(bad=0\land 0\le off\le1020\land
        4\le1024-off\le1024\bigr)
\]
\end{verified}

At the end of Week~10, this was stronger than the earlier
\(bad\in\{0,1\}\) control theorem, but did not state that the entire suffix
returned by the actual procedure equaled
\(
\operatorname{TraceBytes}
  (\operatorname{SymbolList}(symbols_0))
\).
The inner phase then used a conservative \(0..4\) range.
Procedure-level composition of the actual outer iteration's one-symbol state
transition and final four stores remained open.

\subsection{Stage 11: global trace closure of the actual encoder}

Week~11 added \WeekElevenTargetCount\ encoder targets, composing Week~10's
local transitions into the entire successful output
\cite{week11-report,rans-encoder-invariant}.
The central invariant remembers both the initial array \(enc_0\) and the
already processed pure suffix, rather than an arbitrary inner-loop entry array.

\begin{verified}[Tail-aware invariant and generated arithmetic]
The proposition \identifier{encoder_inner_tail_inv} states that the actual
inner loop preserves
\[
  \operatorname{InnerWrittenSuffix}(x_0,s,k)
  \mathbin{+\!\!+}\operatorname{EncodeTrace}(tail).bytes
\]
after the current cursor.

Lemma \theoremname{encoder_inner_tail_exit_exact} recovers
\(k=\operatorname{NormalizationLen}(x_0,s)\le2\) when the actual guard exits.
Lemma \theoremname{generated_loaded_nested_word_update} transports actual
table loads, reciprocal high products, shifts, complements, and biases
to the pure fast step.
\end{verified}

\begin{verified}[Encoder trace-closure theorem]
\theoremname{actual_rans_encode_trace_closure} directly opens actual generated
\identifier{RansEncodeTarget.M._rans_encode}. If it returns, it either fails
or, on success, establishes all of
\[
\begin{gathered}
  bad'=0,\qquad 0\le off'\le1020,\qquad
  4\le1024-off'\le1024,\\
  \operatorname{SegmentMatches}
    \bigl(enc',off',
      \operatorname{TraceBytes}(\operatorname{SymbolList}(symbols_0))\bigr),\\
  off'+\left|
    \operatorname{TraceBytes}(\operatorname{SymbolList}(symbols_0))
  \right|=1024,\qquad
  \operatorname{PrefixFrame}(enc_0,enc',off')
\end{gathered}
\]
\theoremname{actual_rans_encode_trace_refinement} is the corollary expanding
the same conclusion on the parent refinement surface.
\end{verified}

The theorem also composes the fact that the final four generated stores prepend
the serialized state to the accumulated normalization suffix.
Thus \claimid{OBL-RANS-ACTUAL-INNER-NORMALIZE},
\claimid{OBL-RANS-FINAL-SERIALIZATION}, and
\claimid{OBL-RANS-ENCODE-REFINEMENT} are each \statusproved\ within their stated
scopes. This remains Hoare partial correctness: it gives neither termination
of an actual call nor existence of an execution with \(bad'=0\).
The theorem alone therefore cannot establish success for a fixed input.
At Week~11, \claimid{OBL-RANS-ACTUAL-SUCCESS-WITNESS} was \statuspartial;
Week~15's separate all-6 theorem excludes failure in returning fixed-input
executions.

\subsection{Stage 12: exact-trace semantics of the actual decoder}

Week~12 added \WeekTwelveTargetCount\ decoder targets, installing pure-trace
coordinates into the actual generated decoder's two nested loops
\cite{week12-report,rans-decoder-invariant}.
The central input predicate \identifier{actual_mode2_decoder_trace_input}
does not accept an arbitrary buffer; it jointly requires the following.

\begin{itemize}
  \item The expected symbol array is a canonical mode~2 symbol stream;
  \item \(encoded\_size=|\operatorname{TraceBytes}
        (\operatorname{SymbolList}(expected))|\) and \(4\le encoded\_size\le1024\);
  \item \(buffer_0[0..encoded\_size)\) agrees pointwise with that exact trace;
  \item Each generated normalization read obtains the corresponding byte of
        the pure trace segment;
  \item Actual state fields are \((count,size,m)=(1024,encoded\_size,13)\);
  \item The actual mode~2 \identifier{symbol_words} and
        \identifier{dsyms_words} are used.
\end{itemize}

\begin{verified}[Decoder trace coordinates and invariant]
\identifier{decoder_state_at}, \identifier{decoder_cursor}, and
\identifier{decoder_segment_at} turn the pure \(xs/cuts/bytes\) trace's state,
cursor, and normalization segment into per-iteration coordinates.
The proposition \identifier{decoder_outer_trace_live} jointly preserves
\[
  x=\operatorname{DecoderStateAt}(i),\qquad
  off=\operatorname{DecoderCursor}(i)
\]
with equality of the decoded prefix and an untouched-tail frame.
The inner invariant replays exactly 0/1/2 bytes of the corresponding segment,
restoring the next pure state.
\end{verified}

\begin{verified}[Actual decoder exact-trace theorem]
\theoremname{actual_rans_decode_trace_refinement} directly opens actual generated
\identifier{RansDecodeTarget.M._rans_decode} and shows, if it returns,
\[
\begin{gathered}
  bad'=0,\qquad off'=encoded\_size,\\
  \forall i,\;0\le i<1024\Longrightarrow decoded'_i=expected_i,\\
  \forall i,\;1024\le i<2048\Longrightarrow decoded'_i=decoded_{0,i}
\end{gathered}
\]
It also derives the internal exit state \(x=2^{23}\), passing the actual final
reject test. Because the generated ABI does not return \(x\), however, only
returned arrays and state fields are exposed in the postcondition.
\end{verified}

Thus \claimid{OBL-RANS-DEC-NORMALIZE} and
\claimid{OBL-RANS-DECODE-REFINEMENT} are \statusproved\ as exact-trace Hoare
partial correctness.
At the end of Week~12, no sequential-harness theorem transported the encoder's
\identifier{segment_matches} and actual copying's \identifier{slice_eq}
into this input predicate. The next Week~13 stage closes precisely that edge.

\subsection{Stage 13: the actual encoder--copy--decoder core inverse}

Week~13 added \WeekThirteenTargetCount\ targets and composed the three calls
of actual \identifier{Mode2RansActualHarness.run} in order
\cite{week13-report,rans-core-composition}.
The harness and control theorem are in \sourcepath{Mode2RansActualInverse.ec};
semantic closure is separated into
\sourcepath{Mode2RansCoreCompositionBridge.ec} and
\sourcepath{Mode2RansCoreActualInverse.ec}.
The new bridge provides these transport lemmas.

\begin{itemize}
  \item \theoremname{copied_suffix_is_exact_trace}: obtains an exact trace at
        offset~0 of the copied buffer from encoder
        \identifier{segment_matches} and actual-copy \identifier{slice_eq};
  \item \theoremname{trace_bytes_trace_segment_nth}: decomposes global
        \identifier{trace_bytes} coordinates into per-symbol local-segment
        and decoder-cursor coordinates;
  \item \theoremname{segment_matches_implies_exact_decoder_segment_input}:
        obtains every pointwise decoder-byte-read premise from the global trace
        segment; \theoremname{segment_matches_implies_decoder_word_reads}
        transports these into actual W64 reads;
  \item \theoremname{actual_decoder_input_from_copied_trace} and
        \theoremname{configured_decoder_state_fields}: turn the copied trace
        and the three actual state stores, respectively, into components of the
        decoder input relation;
  \item \theoremname{actual_decoder_input_from_configured_trace}: constructs
        the entire Week~12 decoder input predicate from the copied trace,
        canonical symbol stream, size bound, and three actual state stores;
  \item \theoremname{encoder_success_size_word_bridge}: identifies actual
        (W64) subtraction with integer trace size in the no-wrap range.
\end{itemize}

\begin{verified}[\claimid{OBL-RANS-CORE-INVERSE}]
The semantic preconditions of
\theoremname{actual_rans_encode_copy_decode_inverse} are only harness-argument
binding and \identifier{mode2_hbz_symbol_stream}.
It does not assume encoder success, copied-trace equality, decoder execution,
or decoder success. If it returns, one of two branches holds:
\[
\begin{aligned}
\mathsf{failure}:\quad
  &encoderBad\ne0,\quad \neg decoderRan,\quad decoderOff=0,\quad decoderBad=1,\\
  &decoded'=decoded_0,\quad decoderState'=decoderState_0;\\[0.3em]
\mathsf{success}:\quad
  &encoderBad=0,\quad decoderRan,\quad (count,size,m)=(1024,encodedSize,13),\\
  &decoderBad=0,\quad decoderOff=encodedSize,\\
  &decoded'[0..1024)=expected[0..1024),\\
  &decoded'[1024..2048)=decoded_0[1024..2048).
\end{aligned}
\]
The success branch also preserves the actual encoder-suffix relation and initial
prefix frame. The core inverse of the three actual generated procedures
therefore compiles as success-conditioned partial correctness.
\end{verified}

\theoremname{core_composition_preconditions_satisfiable} uses an all-zero
canonical symbol array to show only nonvacuity of the premises.
Because the conclusion is a failure/success disjunction, neither reachability
of the actual success branch nor encoder/decoder termination follows.
Thus \claimid{OBL-RANS-ACTUAL-SUCCESS-WITNESS} remained \statuspartial\
at Week~13. This conclusion did not automatically lift to production full-HBZ
wrappers; the separate lift was closed in Week~14 and fixed all-6 failure
exclusion in Week~15.

\subsection{Stage 14: the production full-HBZ wrapper boundary}

Week~14 added \WeekFourteenTargetCount\ targets connecting focused full-HBZ
procedures to the actual SignaturePack/Unpack extraction boundary
\cite{week14-report,hbz-full-wrapper-composition}.
The proof chain has four layers.

\begin{enumerate}
  \item \sourcepath{Mode2HbzInternalBoundaries.ec} exactly connects prepare,
        rANS encode/decode, and apply calls inside the focused full wrappers
        to existing actual-procedure theorems.
  \item \theoremname{actual_encode_hb_z1_full_mode2_trace} preserves the
        failure/success branches of the actual full encoder on canonical HBZ
        coefficients. On success it provides a prepared-symbol witness,
        exact trace size/segment, and suffix frame.
  \item \theoremname{actual_decode_hb_z1_full_mode2_inverse} shows that,
        given that exact trace and prepared-symbol relation, the actual full
        decoder returns \(bad=0\), the original 1024-coefficient prefix,
        and the coefficient-tail frame.
  \item \theoremname{actual_hbz_full_encode_decode_inverse_mode2} composes
        the focused encode/decode harness, and
        \theoremname{signature_pack_unpack_hbz_full_actual_exact} matches
        it exactly to the production SignaturePack/Unpack harness, including
        the result tuple and \identifier{decoder_ran}.
\end{enumerate}

\begin{verified}[\claimid{OBL-SIG-HBZ-ENCODE-DECODE}]
The only semantic preconditions of the Hoare corollary
\theoremname{signature_pack_unpack_hbz_full_inverse_mode2} are argument binding
and \identifier{canonical_hbz_mode2}. If execution returns, this disjunction
holds:
\[
\begin{aligned}
\mathsf{failure}:\quad
  &size=0,\quad \neg decoderRan,\quad encoded=out_0,\\
  &decoded=decoded_0,\quad bad=bad_0;\\[0.3em]
\mathsf{success}:\quad
  &size\ne0,\quad 4\le size\le1024,\quad decoderRan,\quad bad=0,\\
  &decoded[0..1024)=hbz_0[0..1024),\\
  &\operatorname{CoeffTailFrame}(decoded_0,decoded),\quad
    \operatorname{SuffixFrame}(out_0,encoded),
\end{aligned}
\]
and the success branch has a prepared-symbol witness with
\identifier{segment_matches} for its \identifier{mode2_trace_bytes}.
Production-wrapper composition is therefore \statusproved\ as
success-conditioned partial correctness.
\end{verified}

The theorem does not assume encoder success; it preserves the actual failure
branch. The Week~14 theorem itself, however, proves neither failure exclusion
for a returning execution on canonical all-zero HBZ entering the actual
encoder's nonzero-size branch, nor encoder/decoder termination.
The fact that prepare produces an all-6 symbol stream and the actual encoder's
success are distinct propositions.
The former is composed with Week~15's separate budget/failure-cause theorem;
termination remains separate.

\subsection{Stage 15: fixed all-6 actual-success postconditions}

Week~15 added \WeekFifteenTargetCount\ targets, closing
\claimid{OBL-RANS-ACTUAL-SUCCESS-WITNESS} as fixed-input Hoare partial
correctness \cite{week15-report,rans-actual-success-witness}.
The first target, \sourcepath{Mode2RansAllSixBudget.ec}, proves the following
pure/prepare chain.

\begin{itemize}
  \item \theoremname{zero_hbz_get32}, \theoremname{zero_hbz_canonical}, and
        \theoremname{actual_prepare_zero_hbz_all_six} establish \(bad=0\)
        and a 1024-element all-6 prefix in the actual prepare return;
  \item \theoremname{symbol6_normalization_len_le1} shows that each symbol~6
        step emits at most one normalization byte;
  \item \theoremname{all_six_first_four_no_normalization} and
        \theoremname{all_six_normalization_budget} combine zero-byte emission
        in the first four steps with a total upper bound of 1020 bytes;
  \item \theoremname{all_six_trace_fits_mode2} adds the final four state bytes,
        placing total trace size in \([4,1024]\).
\end{itemize}

The existing encoder closure is strengthened by
\theoremname{actual_rans_encode_failure_trace_cause}, which concludes that
nonzero \(bad\) implies a normalization trace exceeding 1020 bytes.
Contradiction with the all-6 bound gives the following actual postconditions.

\begin{verified}[\claimid{OBL-RANS-ACTUAL-SUCCESS-WITNESS}]
\theoremname{actual_rans_encode_all_six_success} establishes \(bad=0\),
\(0\le off\le1020\), \(4\le1024-off\le1024\), an exact trace segment,
and a prefix frame for every returning execution of actual
\identifier{_rans_encode}.
\theoremname{actual_encode_hb_z1_full_zero_success} and
\theoremname{signature_pack_hbz_zero_success_mode2} lift this to the
focused/production full encoder.
\theoremname{signature_pack_unpack_hbz_zero_success_mode2} concludes the
production zero round trip, \identifier{decoder_ran}, decoder \(bad=0\),
and coefficient/encoded-tail frames.
\end{verified}

This proof does not require success as a precondition, but remains partial
correctness in Hoare logic. It includes no assertion that an actual fixed-input
execution terminates or exists, no losslessness or probability-one success,
and no \identifier{phoare} theorem.
Both “fixed all-6 Hoare success is \statusproved” and “fixed-input termination
is unproved” must therefore be retained.

\subsection{Stage 16: the frozen actual KeyGen snapshot and KG--NTT stopping boundary}

Week~16 added \WeekSixteenTargetCount\ MINCORE targets: three for KeyGen,
one for Sign, and four for Verify
\cite{week16-mincore-plan,week16-kg-report,week16-kg-ntt-mul-report,
week16-sign-report,week16-verify-report}.
In the first pure file,
\sourcepath{Mode2KeygenSnapshot}%
\allowbreak\sourcepath{Algebra.ec},
\theoremname{snapshot_residue_exact_low_high},
\theoremname{raw_sum_raw_residue_congruent_mod_2q}, and
\theoremname{snapshot_expression_from_product_congruent_mod_2q}
prove low/high decomposition of canonical residues and the doubled congruence.

File \sourcepath{Mode2KeygenCoreEquation.ec} defines the transparent harness
\begin{quote}
\identifier{ActualM23MatrixFinalizeSnapshot.run}.
\end{quote}
It snapshots the return of actual
\identifier{_kp_m23_matrix(2,3)}, then calls actual
\identifier{_keypair_finalize_m23(512)}.

\begin{verified}[Actual snapshot/finalizer boundary]
\theoremname{actual_m23_matrix_finalize_snapshot} composes the matrix-output/NTT
representation, exact finalization, and tail frame.
The strongest theorem,
\theoremname{actual_m23_matrix_finalize_semantic_snapshot}, adds
\theoremname{finalize_semantic_output_low_high_decomposition}, HAETAE finalizer
semantics, and \theoremname{finalize_semantic_output_snapshot_mod2q_zero},
yielding a KG-2-like decomposition, adjusted \(s_2\), and a snapshot-only
zero congruence mod-\(2q\) at every active coefficient.
\end{verified}

The third target, \sourcepath{Mode2KeygenNttMulBridge.ec}, fixes only the
last sound rewrite in the existing NTT chain, without duplicating actual loops.
Under the transformed-secret representation,
\theoremname{output_row_from_mode2_ntt_words} proves that
\identifier{output_row} equals exactly
\[
  \operatorname{array256\_mont}
  \bigl(\operatorname{full\_invntt}
    (\operatorname{pointwise\_row\_words}(m,v,row))\bigr)
\]
\theoremname{output_row_repr_from_mode2_ntt_words} transports this equality
to the existing \identifier{wide_slice_repr_bound} contract.
\theoremname{actual_m23_matrix_snapshot_rows_explicit} reuses the strongest
snapshot theorem by consequence, reducing both active rows to that
full-invNTT/pointwise-word representation.
None of these theorems assumes a native \(R_q\) product, security-model product,
or the desired KeyGen equation.

\begin{openobligation}[Frozen \claimid{OBL-MINCORE-KEYGEN}: \identifier{STOP-KG-NTT}]
At the frozen-paper stage, the final NTT-word rewrite compiled, but the
odd-root orthogonality/full-NTT convolution theorem converting actual
pointwise/inverse-NTT representation to native \(R_q\) negacyclic multiplication
was missing, as was the adapter transporting \identifier{Rq.poly} to the
security model's integer-list product. The row product
\(\sum_j A_{\mathrm{gen},ij}s_{\mathrm{gen},j}\) was therefore unavailable.
Faithful KG-1, KG-3, complete KG-4, and the paper equation
\(A s=qj\pmod{2q}\) were not proved.
\end{openobligation}

The KG-NTT-MUL audit identified these two missing leaves and ended at
\identifier{STOP-KG-NTT}, not \identifier{GO-KG}.
KeyGen was frozen at the KG-2/finalization boundary.
Sampling, retries, packing, the public API, KeyGen termination, and security
lie outside the conclusions of this actual two-call snapshot theorem.
Sign was subsequently frozen at \identifier{STOP-SIGN-CHAL-MODE2}, and
Verify at \identifier{STOP-VERIFY-MATRIX-CRT}
\cite{week16-sign-report,week16-verify-report,paper-freeze-artifacts}.

\subsection{Stage 17: post-freeze NTT/Rq transport and the actual avec/qj boundary}

Post-freeze work closed the first algebraic gap above without expanding the
frozen 82-target manifest.
In \sourcepath{NTTFullSpectralAction.ec},
\theoremname{rq_mul_coeff_foldr_to_bigi} and
\theoremname{full_ntt_montgomery_spectral_action} connect the
implementation-style negacyclic coefficient sum to full-NTT spectral action.
\theoremname{full_ntt_montgomery_row_product} composes this into a row product
for arbitrary column counts; \sourcepath{NTTMode2RowSpecializations.ec}
provides pure specializations for KeyGen's three columns and Verify's four.
This Verify result is not a theorem about actual
\identifier{_polyvec_ntt}/\identifier{_polymat_pointwise_acc}/
\identifier{_polyvec_invntt} procedures and does not close frozen
\identifier{STOP-VERIFY-MATRIX-CRT}.

\theoremname{rq_poly_add_repr}, \theoremname{rq_poly_mul_repr}, and
\theoremname{rq_poly_dot3_repr} in \sourcepath{RqHAETAEBridge.ec}, together
with \theoremname{keygen_mode2_row_product},
\theoremname{mode2_row_product_haetae_dot}, and
\theoremname{actual_m23_matrix_snapshot_rows_row_product} in
\sourcepath{Mode2KeygenNttMulBridge.ec}, transport the native \(R_q\)
three-term row product to the HAETAE \identifier{poly_dot} representation
\cite{post-freeze-kg-rq-haetae-report}.

\begin{verified}[Deterministic boundary of actual avec and the paper's \(a,qj\)]
\sourcepath{KgActualAvecQjSemantics.ec} fixes mode~2 row nonces 515/516,
\([3+row,2]\) after the 32-byte seed, SHAKE128 domain/padding 31/128,
168-byte blocks, little-endian 16-bit parsing, and \(<64513\) rejection.
\theoremname{mode2_uniform_stream_paper_a_flat} interprets the two
256-coefficient rows of the terminating actual sampler as the paper's \(a\).
The paper's \(j=(1,0)^T\) is separated through
\theoremname{paper_qj_row0_coeff0},
\theoremname{paper_qj_row0_other}, and
\theoremname{paper_qj_remaining_row}.
\theoremname{checked_mode2_parent_m23_finalize_actual_avec_qj} combines
current actual-sampler, matrix-helper, and finalizer results into one Hoare
boundary. This is partial correctness of terminating executions, not a
distribution, uniformity, losslessness, or termination claim.
\end{verified}

\begin{verified}[Faithful augmented coefficient-product composition]
The latest unmanifested standalone file
\sourcepath{KgFaithfulAugmentedModel.ec} instantiates a width-6 carrier with
actual arrays through
\identifier{faithful_augmented_matrix_coeff},
\identifier{faithful_augmented_secret_coeff},
\identifier{actual_faithful_augmented_matrix}, and
\identifier{actual_faithful_augmented_secret}.
\identifier{faithful_middle_matrix_poly} and
\identifier{faithful_middle_secret_poly} reconstruct the carrier's middle
three columns as HAETAE polynomials, so there is no free generated-block
argument detached from the displayed matrix/secret.
It then composes
\theoremname{actual_snapshot_paper_qj_to_faithful_augmented},
\theoremname{generated_product_coeff_from_rq_repr},
\theoremname{faithful_augmented_As_qj_from_generated_repr},
\theoremname{rq_as_haetae_poly_repr},
\theoremname{faithful_generated_row_product_haetae_dot},
\theoremname{actual_faithful_augmented_productE},
\theoremname{faithful_augmented_product_qj_from_generated_dot}, and
\theoremname{mode2_m23_facts_generated_row_product}.
The strongest theorem,
\theoremname{checked_mode2_parent_m23_finalize_faithful_augmented_paper_as_qj},
proves congruence mod-\(2q\) between
\identifier{faithful_augmented_matrix_vector_product_coeff} and
\identifier{faithful_qj_coeff} at each active coefficient of the actual
checked sampler/matrix/finalizer parent.
The file passed standalone \identifier{-no-eco} fresh compilation.

The exact endpoint is a structured coefficient product specific to the actual
carrier. \identifier{actual_faithful_augmented_productE} closes the expansion
of the product taking only matrix and secret as inputs.
The middle 2-by-3 block is reconstructed from the specified columns and
computed with HAETAE \identifier{poly_dot}.
The verdict for this work is \FaithfulAugmentedVerdict.
These dedicated mod-\(2q\) block semantics are not a coefficient theorem for
the existing width-4, mod-\(q\) security \identifier{matrix_vec_mul}, nor do
they identify security-side synthetic ExpandVecA with the actual sampler.
\cite{post-freeze-kg-faithful-augmented-report}
\end{verified}

\begin{openobligation}[\CurrentGate: \claimid{OBL-KG-SECURITY-EXPANDVECA-SEMANTICS}]
The existing security-side \identifier{haetae_mode23_qj_vector} is a
synthetic object using \identifier{public_key_seed_polyveck} tag 401.
\theoremname{security_named_qj_is_not_paper_qj} proves that its row~0
coefficient~0 is 401, whereas the paper's \(qj\) has \(q=64513\).
The zero-seed execution also has actual nonce-515 first coefficient 44985.
The security specification must therefore be replaced or refined with actual
SHAKE128/LE16/rejection/nonce semantics, keeping the paper's \(qj\) as a
separate correction, and an actual-array representation theorem must then
be proved. The paper-shaped congruence for the new custom carrier cannot
therefore be promoted to KG-1--KG-4 of the existing security model or
security-level \(A s=qj\pmod{2q}\)
\cite{post-freeze-kg-actual-avec-qj-report}.
\end{openobligation}

\subsection{Exact endpoints of the completed chains}

\begin{longtable}{@{}L{100mm} L{55mm}@{}}
\toprule
Chain & Endpoint \\
\midrule
\endfirsthead
\toprule
Chain & Endpoint \\
\midrule
\endhead
packed key \(\to\) generated raw \(\mu\) \(\to\) position-64 absorb &
\statusproved\ (raw/internal, local) \\
actual KeyGen export, Sign/Verify import, address/frame/control &
\statusproved\ (each named contract) \\
direct equality of stored \(\mu\) in completed actual traces &
\statuspartial \\
canonical signature prefix round trip &
\statusproved\ (1056-byte prefix) \\
HBZ leaf, table, pure rANS step and byte trace &
\statusproved\ (pure/leaf) \\
actual encoder array/list, word, and serialization leaves &
\statusproved\ (Week~10 local propositions) \\
actual encoder inner byte/frame and successful-output size &
\statusproved\ (exact \(0..2\), actual partial correctness) \\
actual encoder entire trace-suffix refinement &
\statusproved\ (success-conditioned partial correctness) \\
actual decoder exact-trace refinement &
\statusproved\ (1024-symbol recovery, exact consumption, tail frame) \\
actual encoder--copy--decoder harness control &
\statusproved\ (calls/branches only) \\
actual encoder--copy--decoder inverse &
\statusproved\ (success-conditioned partial correctness) \\
production full-HBZ wrapper inverse &
\statusproved\ (success-conditioned partial correctness, production boundary) \\
actual encoder success witness &
\statusproved\ (Hoare failure exclusion for fixed all-6 input;
no termination/losslessness) \\
actual KeyGen matrix/finalize snapshot &
\statuspartial\ (KG-2-like decomposition and snapshot congruence proved;
frozen paper ends at \identifier{STOP-KG-NTT} after the final NTT-word rewrite) \\
post-freeze full-NTT/native \(R_q\)/HAETAE row-product transport &
\statusproved\ (named pure/representation theorems outside the frozen 82 targets) \\
post-freeze actual \identifier{avec} and paper \(a,qj\) boundary &
\statuspartial\ (deterministic semantics and separate \(qj\) composition proved;
\CurrentGate\ at the security-object mismatch) \\
faithful augmented coefficient-product/transport &
\statusproved\ (standalone, outside the frozen 82 targets;
actual-carrier-specific product, not generic security-model
\identifier{matrix_vec_mul} equivalence) \\
actual accepted Sign core &
\statuspartial\ (call/branch control only; \identifier{STOP-SIGN-CHAL-MODE2}) \\
actual Verify core &
\statuspartial\ (helper-local word/control only; frozen \FrozenPaperGate) \\
\(h\) suffix codec &
\statusdeferred\ (Week~16 MINCORE KeyGen scope) \\
full signature codec and Sign--Verify correctness &
\statusblocked \\
implementation-level functional correctness and security &
\statusspecified / \statuspartial \\
\bottomrule
\end{longtable}
